%% shortname: Expected Utility, Money Lotteries and Risk Aversion
\documentclass[11pt]{article}
\usepackage{microstyle}

\begin{document}

\lessonheader{16}{Chapter 6 (continued): the expected utility theorem, money lotteries and risk aversion}{}


\section{Recap: preferences over lotteries}

We introduced preferences $\succsim$ over lotteries.
\begin{flatlist}
  \item Complete $+$ transitive $+$ continuity $\Rightarrow$ preferences can be
        represented with a utility function $U:\mathcal{L}\to\R$.
  \item The independence axiom adds structure: it implies that indifference curves
        are straight lines (or hyperplanes) and that they shift in a parallel
        fashion.
\end{flatlist}


\section{The expected utility theorem}

The central result is expected utility.

\topic{Theorem (expected utility)}
If the axioms are satisfied (rational, continuous, independence), then $\succsim$ can
be represented by a utility function with the expected utility property:
\[ U(L)=\sum_{n=1}^{N}p_n u_n, \qquad L=(p_1,\dots,p_N), \]
where $p_1,\dots,p_N$ are the probabilities of the lottery and $u_n$ is the utility if
outcome $n$ occurs.

\runin{Sanity check.} Indifference curves are straight lines because $U(L)$ is
linear in the probabilities, and they also shift in a parallel way. An indifference
curve is
\[ U(L)=c \quad\Longleftrightarrow\quad p\cdot u=c, \qquad
   (p_1,\dots,p_N)\cdot(u_1,\dots,u_N)=c, \]
a hyperplane in $p$ with normal vector $u=(u_1,\dots,u_N)$. Given $u_1,\dots,u_N$, the
hyperplane shifts in parallel when $c$ changes.

\runin{Note.} It is not claimed that the $u_n$'s are unique. The same preferences are
represented after a positive affine transformation: change $u_n$ to
\[ \alpha+\beta u_n, \qquad \beta>0, \qquad \forall\, n=1,\dots,N. \]
Then $\sum_n p_n(\alpha+\beta u_n)=\alpha+\beta\sum_n p_n u_n$ (the $p_n$ sum to $1$),
so the ranking of lotteries is unchanged.


\section{Money lotteries}

Consider now lotteries where the outcome is money.

It is convenient to think of money as continuous, but we then need to make some
assumptions, because we can no longer express the lottery as a list of
probabilities. How do we represent a lottery over a continuous variable?

\topic{Cumulative distribution functions}
We use cumulative distribution functions (CDF),
\[ F:\R\to[0,1], \qquad F(x)=\Pr(\tilde x\le x), \]
the probability that the outcome is $x$ or less.
\begin{flatlist}
  \item Clearly, $F(x)$ is non-decreasing.
  \item Not necessarily continuous, but always right-continuous.
  \item Any real random variable has a CDF.
\end{flatlist}

\begin{center}
\begin{tikzpicture}[scale=1.05]
  \ecaxes{4.6}{3.0}
  \node[ecnode,right] at (4.6,0) {$x$};
  \node[ecnode,above] at (0,3.0) {Prob.};
  \draw[dashed] (0,2.4) -- (4.3,2.4);
  \node[ecnode,left] at (0,2.4) {$1$};
  \draw[budget] (0,0.2) .. controls (1.2,0.5) and (2.6,1.4) .. (4,2.4) -- (4.3,2.4);
  \node[ecnode,below right] at (3.3,1.85) {$F(x)$};
  % x' = 1.2, x'' = 2.7
  \draw[dashed] (1.2,0) node[ecnode,below] {$x'$} -- (1.2,0.52) -- (0,0.52)
                node[ecnode,left] {$F(x')$};
  \draw[dashed] (2.7,0) node[ecnode,below] {$x''$} -- (2.7,1.45) -- (0,1.45)
                node[ecnode,left] {$F(x'')$};
  \node[ecnode,below] at (4,0) {$\bar x$};
  \draw[dashed] (4,0) -- (4,2.4);
\end{tikzpicture}
\end{center}
\begin{align*}
  F(x'') &= \Pr(\tilde x\le x''),\\
  F(x')  &= \Pr(\tilde x\le x'),\\
  F(x'')-F(x') &= \Pr\bigl(\tilde x\in(x',x'']\bigr).
\end{align*}
As $x'\to x''$ (from below), $F(x'')-F(x')\to\Pr(\tilde x=x'')$; so the difference goes
to $0$ if there is zero probability that $\tilde x=x''$.

\begin{center}
\begin{tikzpicture}[scale=1.05]
  \ecaxes{4.6}{3.0}
  \node[ecnode,right] at (4.6,0) {$x$};
  \node[ecnode,below left] at (0,0) {$0$};
  \draw[dashed] (0,2.6) -- (4.3,2.6);
  \node[ecnode,left] at (0,2.6) {$1$};
  % continuous part up to x' = 1.3, jump at x', rise to x'' = 2.0, flat to x''' = 3.2, rise to 1
  \draw[budget] (0,0) .. controls (0.3,0.7) and (1.0,0.7) .. (1.3,1.0);
  \draw[budget] (1.3,1.55) .. controls (1.6,1.7) and (1.8,1.95) .. (2.0,2.0) -- (3.2,2.0)
                .. controls (3.6,2.1) and (3.9,2.5) .. (4.1,2.6) -- (4.3,2.6);
  \draw[fill=white] (1.3,1.0) circle (1.6pt);
  \node[bundle] at (1.3,1.55) {};
  \draw[dashed] (1.3,0) node[ecnode,below] {$x'$} -- (1.3,1.0);
  \draw[dashed] (0,1.0) node[ecnode,left] {$\Pr(\tilde x<x')$} -- (1.3,1.0);
  \draw[dashed] (0,1.55) node[ecnode,left] {$F(x')$} -- (1.3,1.55);
  \draw[dashed] (0,2.0) node[ecnode,left] {$F(x'')=F(x''')$} -- (2.0,2.0);
  \draw[dashed] (2.0,0) node[ecnode,below] {$x''$} -- (2.0,2.0);
  \draw[dashed] (3.2,0) node[ecnode,below] {$x'''$} -- (3.2,2.0);
  \node[ecnode,align=left,anchor=west] at (5.0,1.5)
       {positive probability of $x'$:\\ $\Pr(\tilde x=x')=F(x')-\Pr(\tilde x<x')>0$;\\[2pt]
        zero probability of $x\in(x'',x''')$:\\ $F$ is flat there.};
\end{tikzpicture}
\end{center}

\topic{Expected utility for money lotteries}
Turning to expected utility, define $u(x)$ as the Bernoulli utility of $x$. Assume $x$
is bounded below, and assume that $u(x)$ is increasing, continuous and bounded. We can
now calculate utility:
\[ U(F)=\int u(x)\,dF(x). \]
If $F(x)$ has a density, $F'(x)=f(x)$, then
\[ U(F)=\int u(x)f(x)\,dx. \]


\section{Risk aversion}

Preferences depend on $u(x)$.

\topic{Definition}
The agent is risk averse if $u(x)$ is concave. Such an agent prefers getting the
expected value of $F$, $\int x\,dF(x)$, for sure over taking the risky lottery $F$.
This follows from Jensen's inequality:
\[ \int u(x)\,dF(x)\ \le\ u\Bigl(\int x\,dF(x)\Bigr). \]

\begin{center}
\begin{tikzpicture}[scale=1.1]
  % u(x) = 1.5 sqrt(x); lottery x_L = 0.25, x_H = 3.75 with prob. 1/2 each
  \ecaxes{4.6}{3.4}
  \node[ecnode,right] at (4.6,0) {$x$};
  \draw[budget,domain=0:4.2,samples=60] plot (\x,{1.5*sqrt(\x)});
  \node[ecnode,right] at (4.2,3.07) {$u(x)$};
  \draw (0.25,0.75) -- (3.75,2.905);                             % chord
  \node[bundle] at (0.25,0.75) {}; \node[bundle] at (3.75,2.905) {};
  \node[bundle] at (2,2.121) {};                                % u(E x)
  \node[bundle] at (2,1.827) {};                                % E u(x)
  \draw[dashed] (0.25,0) node[ecnode,below] {$x_L$} -- (0.25,0.75);
  \draw[dashed] (3.75,0) node[ecnode,below] {$x_H$} -- (3.75,2.905);
  \draw[dashed] (2,0) node[ecnode,below=11pt] {$E(x)$} -- (2,2.121);
  \draw[dashed] (1.484,0) node[ecnode,below] {CE} -- (1.484,1.827);
  \draw[dashed] (0,1.827) -- (2,1.827);
  \draw[dashed] (0,2.121) -- (2,2.121);
  \node[ecnode,left] at (0,2.18) {$u(E(x))$};
  \node[ecnode,left] at (0,1.78) {$\int u\,dF$};
\end{tikzpicture}
\end{center}
As the graph shows for risk aversion: the expected utility of the lottery (on the
chord) lies below the utility of its expected value (on the curve).

\topic{Concavity and attitude towards risk}
\begin{flatlist}
  \item $u''(x)<0\ \ \forall x \ \Rightarrow$ risk averse;
  \item $u''(x)=0\ \ \forall x \ \Rightarrow$ risk neutral;
  \item $u''(x)>0\ \ \forall x \ \Rightarrow$ risk seeking (risk loving).
\end{flatlist}


\section{Measuring the degree of risk aversion}

\topic{Certainty equivalent}
For a given $F$, the certainty equivalent $c(F,u)$ (CE) is the amount of money that
makes you indifferent between that amount of money for sure and taking the lottery:
\[ u(\mathrm{CE})=\int u(x)\,dF(x)\ \le\ u\Bigl(\int x\,dF(x)\Bigr)
   \quad\Longleftrightarrow\quad \mathrm{CE}\ \le\ \int x\,dF(x), \]
for a risk-averse agent (the equivalence uses that $u$ is increasing).

\topic{The Arrow--Pratt measure of absolute risk aversion}
Given $u(x)$, the CE is specific to the lottery. We want a measure that does not
depend on $F$.

\runin{Idea.} Measure ``how concave'' $u(x)$ is:
\[ r_A(x\mid u)=-\frac{u''(x)}{u'(x)}\ \ge\ 0 \quad\text{if } u''(x)\le 0,
   \ \text{i.e.\ risk averse.} \]

\runin{Why divide by $u'(x)$?} Because $u(x)$ is not unique: it is only unique up to a
positive affine transformation. Replacing $u(x)$ with $\alpha+\beta u(x)$,
$\beta>0$,
\[ r_A=-\frac{\beta u''(x)}{\beta u'(x)}=-\frac{u''(x)}{u'(x)}, \]
so the measure does not change, whereas $u''(x)$ alone would be scaled by $\beta$.

\runin{Comparing agents.} The higher $r_A$ is, the more risk averse the agent is. If
\[ r_A(x\mid u_1)\ \ge\ r_A(x\mid u_2) \qquad \forall x, \]
then $u_1(x)$ is an increasing and concave transformation of $u_2(x)$:
\[ u_1(x)=T\bigl(u_2(x)\bigr), \qquad T \text{ increasing and concave}. \]
Similarly, $u_1$ has a smaller CE than $u_2$: $c(F,u_1)\le c(F,u_2)$ for every $F$.


\begin{transcriptnotes}
  \item The export has 5 pages; page 1 is the blank notebook cover and is skipped. No
        date or title is written on the pages, so the header date is left empty; the
        subtitle is \textbf{added} (MWG Chapter 6, Sections 6.B--6.C).
  \item Recap: ``Complete + transitive - continuity'' read as complete $+$ transitive
        $+$ continuity. The sentence ``The independence axiom adds structure \ldots
        parallel fashion'' appears twice on page 2 (before and after the theorem);
        kept once. ``indiff'' $\to$ \emph{indifference curves}.
  \item Theorem: ``then $\succsim$ by presented by utility function with the expected
        utility prob.'' $\to$ \emph{can be represented by a utility function with the
        expected utility property}. The page writes $U(L)=\sum_{i=1}^N P_i u_i$ with
        an arrow from $U(L)$ to ``$P_1\cdots P_2$'' (read $p_1,\dots,p_N$) and
        ``utility is outcome occurs'' (read \emph{utility if outcome $n$ occurs});
        typeset with index $n$ as in Lesson 15. ``(rational, continuous,
        independence)'' is \textbf{added}.
  \item Sanity check: ``hyperplane that shift when $C$ changes'' --- the level $c$
        (lower case, not the outcome set $C$). The normal-vector remark is
        \textbf{added}. Note: ``preferences are presented under affine
        transformation'' $\to$ \emph{the same preferences are represented after a
        positive affine transformation}; the one-line check
        $\sum_n p_n(\alpha+\beta u_n)=\alpha+\beta U(L)$ is \textbf{added}.
  \item CDF: ``outcome is $X$ on less'' $\to$ \emph{$x$ or less}; ``all wants
        righ-continuous'' $\to$ \emph{always right-continuous}; ``Ant Real random
        Var. has avilible CDF'' $\to$ \emph{any real random variable has a CDF}. The
        random variable is written as a script $\mathcal{X}$; typeset $\tilde x$.
        ``$\Pr(\tilde x\in(x',x''])$'' kept. \textbf{Corrected:} the page says ``as
        $x'\to x''$, $F(x'')-F(x')\to 0$ if prob zero that $\tilde x=x''$''; the
        limit is $\Pr(\tilde x=x'')$ in general, which is $0$ exactly when there is no
        mass at $x''$ --- both statements are now given.
  \item Second CDF sketch: redrawn to scale. The page's labels $F(x'')=F(x''')$,
        $F(x')$ and ``$\Pr[\tilde x<x']$'' (at the open circle) are kept; the margin
        note ``positive prob of $x'$, zero prob of $x\in(x'',x''')$'' is kept and the
        jump size $F(x')-\Pr(\tilde x<x')$ is \textbf{added}.
  \item Expected utility: ``benooli'' $\to$ \emph{Bernoulli}; ``caunculate'' $\to$
        \emph{calculate}. ``Assuming $x$ is boun DEN belong'' is read as \emph{assume
        $x$ is bounded below} (\textbf{unclear}). \textbf{Corrected:} the page writes
        $U(X)=\int u(x)\,dF(x)$; the utility is a function of the distribution, so it
        is typeset $U(F)$, matching the next line ``Then $U(F)=\int u(x)f(x)\,dx$''.
  \item Risk aversion: \textbf{Corrected:} ``prefers getting expected utility in $F$
        over taking the risky lottery'' $\to$ the \emph{expected value} of $F$,
        $\int x\,dF(x)$, for sure. ``Jensen Equality'' $\to$ \emph{Jensen's
        inequality} (the displayed formula is an inequality). The concave-$u$ sketch
        (with $x_L$, $x_H$, $E(x)$, CE) is redrawn to scale with $u=1.5\sqrt{x}$ and a
        $\tfrac12$--$\tfrac12$ lottery on $x_L=0.25$, $x_H=3.75$ ($E(x)=2$, CE $\approx1.48$).
  \item The boxed note ``Add back clause: $u''<0$ / $=0$ / $>0$ $\forall x$ $\Rightarrow$
        risk averse / neutral / seeking'' (page 5) is inserted in the risk-aversion
        section, where it belongs. Strict concavity gives \emph{strict} risk
        aversion; the definition above (concave, $u''\le0$) is the weak version.
  \item Measuring risk aversion: ``certainty equivlent'' $\to$ \emph{certainty
        equivalent}; ``indifferen'' $\to$ \emph{indifferent}; the MWG notation
        $c(F,u)$ is \textbf{added}. ``ARROR--PRATT'' $\to$ \emph{Arrow--Pratt}. The page
        writes $R_A(X\mid U)$; typeset $r_A(x\mid u)$. ``why $u'(x)$?'' expanded to
        \emph{why divide by $u'(x)$?}; the cut-off right edge of page 4 shows
        $-\beta u''/\beta u'$, completed as $=-u''/u'$.
  \item Comparing agents: ``the higher $R_A$ is, the risk averse the agent is'' $\to$
        \emph{the more risk averse}; ``an $\uparrow$ and concave transformative'' $\to$
        \emph{an increasing and concave transformation}. ``for every $F$'' is
        \textbf{added} (MWG Proposition 6.C.2).
\end{transcriptnotes}

\end{document}
